\section{Conclusions}
This paper examines how preferential college admission policies influence educational outcomes by changing high school students' study incentives, and how students' responses to such incentives are shaped by their perceptions. We leverage a unique experimental setting in Chile—the randomized rollout of the nationwide PACE program---combined with rich administrative and primary survey data. \par

PACE substantially increased selective college admissions and first-year enrollment among disadvantaged students, particularly those at the top of their high school cohorts. However, these gains faded over time, yielding smaller effects on enrollment five years after high school. At the same time, the program generated unintended consequences in the form of reduced pre-college effort and academic achievement.

Our data reveal systematic overconfidence among students about their academic standing and future college success. To examine how beliefs shape behavior and to evaluate counterfactual policies, we estimate a dynamic structural model with subjective expectations that exploits our survey data to relax standard rational-expectations assumptions \citep{manski2004measuring}. The estimated model shows that the short- and longer-term impacts of preferential admissions critically depend on students’ beliefs. Counterfactual simulations where students in both treated and control schools hold rational expectations reveal that, contrary to students’ perceptions, PACE did not offer widespread admission guarantees, but instead modestly increased the returns to effort on average, by making admission attainable. Belief errors were therefore central to the observed effort disincentives.\par




Pairing PACE with an informational intervention that corrects all belief errors, however, would further reduce pre-college effort and weaken enrollment and persistence effects compared to offering PACE alone. This occurs because, regardless of PACE, students’ over-optimism about effort returns and college performance leads them to exert more effort and enroll at higher rates than they would under accurate beliefs, so fully correcting beliefs generates discouragement effects.
\par 

Our results highlight that ignoring students’ subjective beliefs in policy design can generate unintended outcomes. Preferential admissions without carefully tailored informational interventions can discourage effort. We show that interventions emphasizing the role of pre-college effort in college success could mitigate these disincentives while preserving enrollment gains. 
Assessing how improving the college enrollment of disadvantaged students ultimately affects their welfare is beyond the scope of this study.  An important direction for future research is to examine impacts on graduation and labor market outcomes, which are central to evaluating the trade-off between the costs of college attendance and its long-run benefits.  \par 


